平面波量子化已经固定了等时对易关系与能量归一化。截面改变时，这些条件仍须保持；改变的是空间算子的谱及定义域。以下采用与原标量场相容的规范化 $\varphi=\phi_{\rm phys}/\sqrt{\hbar c}$，先在有限模态中建立一般截面的振子，再按谱测度处理连续谱。

\subsection{由正自伴空间算子量子化一般截面}
\label{subsec:ghs-mode-quantization}


先取一个静态的 $(n+1)$ 维洛伦兹时空 $\RR\times\Sigma$，其中时间坐标是具有长度量纲的 $x^0=ct$，$\Sigma$ 是一个紧致、无边界的 $n$ 维黎曼空间，度规为 $h_{ab}$。这里 $n$ 是空间维数，后文讨论欧氏共形场论时则以 $d$ 表示整个欧氏空间维数；在两个问题之间转换时，不应提前把 $n$ 和 $d$ 混用。为使谱从正数开始，先考虑严格正的空间算子
\begin{equation}\label{eq:ghs-bridge-spatial-A}
 A=-\Delta_h+\mu^2,\qquad \mu>0,
\end{equation}
其中 $\mu$ 具有逆长度量纲；若场的物理质量为 $m$，则 $\mu=mc/\hbar$。$\Delta_h=\operatorname{div}_h\nabla_h$ 的本征值取非正的号。对应的无量纲作用量（在合适的场归一化后）为
\begin{equation}\label{eq:ghs-bridge-lorentz-action}
 \mathscr S[\varphi]=\frac12\int\dd x^0\,\dd\mu_h
 \bigl[(\partial_0\varphi)^2-h^{ab}(\partial_a\varphi)(\partial_b\varphi)-\mu^2\varphi^2\bigr].
\end{equation}
物理作用量为 $S=\hbar\mathscr S$；规范化场 $[\varphi]=L^{(1-n)/2}$，因而 $[q_j]=L^{1/2}$、$[p_j]=L^{-1/2}$，$[q_j,p_k]=\ii\delta_{jk}$ 的两边均无量纲。这里 $\dd\mu_h=\sqrt{\det h}\,\dd^ny$。本节选择规范化振子变量；恢复物理量时每个振子的能量是 $\hbar c\omega$，其中 $\omega$ 是逆长度本征频率。以 $x^0$ 作为演化变量使 $c$ 与 $\hbar$ 不会在谱的定义中被暗中设为一。

在无边界紧空间上，$-\Delta_h$ 的非负本征值是离散的。记 $u_j\in L^2(\Sigma,\dd\mu_h)$ 为一组实正交归一的本征函数，重数计入指标 $j$，满足
\begin{equation}\label{eq:ghs-bridge-eigenbasis}
 Au_j=\omega_j^2u_j,\quad \omega_j>0,\qquad
 \int_\Sigma u_j u_k\,\dd\mu_h=\delta_{jk},\qquad
 \sum_j u_j(y)u_j(y')=\delta_h(y,y').
\end{equation}
最后一条是分布意义的完备性关系：$\int\delta_h(y,y')f(y')\dd\mu_h(y')=f(y)$。$u_j$ 是空间函数，$\omega_j$ 是频率而非能量；二者已经足以代替平直空间的 Fourier 波矢。

将场展开为 $\varphi(x^0,y)=\sum_jq_j(x^0)u_j(y)$。同一展开给出
\begin{align*}
 \int_\Sigma(\partial_0\varphi)^2\dd\mu_h
 &=\sum_{j,k}\dot q_j\dot q_k\int_\Sigma u_ju_k\dd\mu_h
 =\sum_j\dot q_j^{\,2},\\
 \int_\Sigma\bigl(|\nabla\varphi|_h^2+\mu^2\varphi^2\bigr)\dd\mu_h
 &=\int_\Sigma\varphi A\varphi\dd\mu_h\\
 &=\sum_{j,k}q_jq_k\omega_k^2\int_\Sigma u_ju_k\dd\mu_h
 =\sum_j\omega_j^2q_j^2.
\end{align*}
第二行的积分分部在无边界紧空间上没有边界项；当 $\Sigma$ 有边界时，必须先选择使 $A$ 自伴的条件，不能直接使用这一等式。于是场作用量完全化成
\begin{equation}\label{eq:ghs-bridge-oscillator-action}
 \mathscr S=\frac12\int\dd x^0\sum_j(\dot q_j^2-\omega_j^2q_j^2).
\end{equation}
这就是“自由场是无限多个振子”的一般数学含义：真正重要的条件是空间算子存在合适的谱分解，而非平面波是否存在。



为避免把无限维问题直接塞进未经定义的乘积，先只保留前 $N$ 个模。共轭动量为 $p_j=\dot q_j$，以逆长度计的哈密顿量为 $H_N=\frac12\sum_{j=1}^{N}(p_j^2+\omega_j^2q_j^2)$。施加 $[q_j,p_k]=\ii\delta_{jk}$，定义
\begin{align*}
 a_j&=\sqrt{\frac{\omega_j}{2}}q_j+\frac{\ii}{\sqrt{2\omega_j}}p_j,\\
 a_j^\dagger&=\sqrt{\frac{\omega_j}{2}}q_j-\frac{\ii}{\sqrt{2\omega_j}}p_j.
\end{align*}
逐项计算交叉对易子得
\begin{align*}
 [a_j,a_k^\dagger]
 &=-\frac{\ii}{2}\sqrt{\frac{\omega_j}{\omega_k}}[q_j,p_k]
 +\frac{\ii}{2}\sqrt{\frac{\omega_k}{\omega_j}}[p_j,q_k]
 =\delta_{jk}.
\end{align*}
反解为 $q_j=(a_j+a_j^\dagger)/\sqrt{2\omega_j}$ 和
$p_j=-\ii\sqrt{\omega_j/2}(a_j-a_j^\dagger)$，代回 $H_N$ 可得
\[
 H_N=\sum_{j=1}^{N}\omega_j\left(a_j^\dagger a_j+\frac12\right).
\]
物理哈密顿量为 $\hbar cH_N$。零点项 $\frac12\sum_j\omega_j$ 在 $N\to\infty$ 时通常发散；后文讨论的正规序哈密顿量只计算相对真空的激发能量，不等于已经解决了一般背景上的真空能量重整化。



有限模真空 $|0_N\rangle$ 由 $a_j|0_N\rangle=0$ 定义。由此制备的任意有限粒子态可以通过产生算符作用于真空得到。无限模极限不再只是形式上的 $N\to\infty$：要得到真正的 Fock 表示，还必须讨论哪些序列有有限范数、不同正频分解之间能否由酉算子联系。这正是后文复结构与 Bogoliubov 理论要回答的问题。


要理解欧氏径向量子化，先看一个振子的真空在坐标表象中是什么。令 $p_j=-\ii\partial_{q_j}$。条件 $a_j\psi_0(q)=0$ 就是
\[
 \left(\sqrt{\frac{\omega_j}{2}}q_j+\frac1{\sqrt{2\omega_j}}\partial_{q_j}\right)\psi_0(q_j)=0,
 \qquad \partial_{q_j}\log\psi_0=-\omega_jq_j.
\]
积分后，利用 $\int_\RR\ee^{-\omega q^2}\dd q=\sqrt{\pi/\omega}$ 归一化，得到
\begin{equation}\label{eq:ghs-bridge-vacuum-one-mode}
 \psi_{0,j}(q_j)=\left(\frac{\omega_j}{\pi}\right)^{1/4}\ee^{-\omega_jq_j^2/2}.
\end{equation}
$N$ 个独立模的真空波函数是这些高斯函数的乘积，其指数为 $-\frac12\sum_{j=1}^{N}\omega_jq_j^2$。将 $q_j=\langle u_j,f\rangle_{L^2}$ 视为边界场 $f(y)$ 的谱系数，就得到形式表达
\begin{equation}\label{eq:ghs-bridge-vacuum-functional}
 \Psi_0[f]\ \propto\
 \exp\left[-\frac12\langle f,A^{1/2}f\rangle_{L^2(\Sigma)}\right].
\end{equation}
这里的 $A^{1/2}$ 定义为 $A^{1/2}u_j=\omega_ju_j$。它一般不是局域微分算子：$A$ 是二阶算子，而其正平方根通常在一点上的取值依赖整个截面上的 $f$。因此真空波泛函的高斯核已经含有整个空间的几何信息。



式 \cref{eq:ghs-bridge-vacuum-functional} 首先是有限模正则化下的高斯泛函。无限维空间没有平移不变的 Lebesgue 测度 $\mathcal Df$；严谨构造应通过圆柱高斯测度或相应的算子代数完成。若 $A^{-1/2}$ 不是迹类算子，典型的真空场不能视为普通 $L^2$ 函数，而是分布。后文的谱核在分离点可按分布方式解释。



波泛函的指数为什么恰好出现平方根，而不是 $A$ 本身？因为 $A$ 给出振子的\emph{频率平方}，但高斯真空由一阶方程 $a_j\psi_0=0$ 决定。取平方根正是把二阶场方程分成正频与负频两部分的一次运算。稍后在一般曲时空上，若不存在一个全局不变的正频划分，真空也就不再有这个唯一的形式。


在初等场论中，我们常在动量空间写传播子 $(p^2+m^2)^{-1}$。如果空间没有平移对称性，动量就不是好量子数；然而“传播子是微分算子的逆”仍然成立。先作 Wick 转动，将 $x^0$ 换为欧氏时间 $\tau$，二次作用量是
\begin{equation}\label{eq:ghs-bridge-Euclidean-action}
 \mathscr I_E=\frac12\int_{\RR}\dd\tau\int_\Sigma\dd\mu_h
 \left[(\partial_\tau\varphi)^2+\varphi A\varphi\right].
\end{equation}
对应欧氏算子为 $L=-\partial_\tau^2+A$。我们要找的两点函数 $G$ 应满足
\begin{equation}\label{eq:ghs-bridge-green-equation}
 L_{(\tau,y)}G(\tau,y;\tau',y')
 =\delta(\tau-\tau')\,\delta_h(y,y').
\end{equation}
这条定义明确了逆算子，也说明传播子依赖边界条件：在整条欧氏时间轴上取两端衰减的解，与有限柱体上固定边界场的解并不相同。

将 $G$ 沿空间本征基展开为 $G=\sum_jg_j(\tau,\tau')u_j(y)u_j(y')$，代入 \cref{eq:ghs-bridge-green-equation} 并使用正交性，可逐模得到
\[
 (-\partial_\tau^2+\omega_j^2)g_j(\tau,\tau')=\delta(\tau-\tau').
\]
当 $\tau\ne\tau'$ 时，齐次方程只允许 $\ee^{\pm\omega_j\tau}$。要求 $|\tau|\to\infty$ 衰减，则解应与 $\ee^{-\omega_j|\tau-\tau'|}$ 成正比。记系数为 $c_j$，将方程跨越 $\tau'$ 的一个小区间积分，得到导数跳跃条件
\[
 -\left[\partial_\tau g_j\right]_{\tau'^-}^{\tau'^+}=1,
 \qquad 2\omega_jc_j=1.
\]
因此任意严格正自伴空间算子的欧氏圆柱传播核为
\begin{equation}\label{eq:ghs-bridge-spectral-G}
 G(\tau,y;\tau',y')
 =\sum_j\frac{\ee^{-\omega_j|\tau-\tau'|}}{2\omega_j}
 u_j(y)u_j(y').
\end{equation}
这就是熟悉的动量积分在一般紧致空间上的替代物。若空间非紧，谱可能包含连续部分，此时必须将求和替换为投影谱测度的积分，不能人为把连续谱离散化。


这一离散本征函数展开还可改写为不依赖球面坐标的谱积分。
对正自伴算子 $A$，谱定理给出投影值测度 $E_A(\lambda)$，并定义 $F(A)=\int_{(0,\infty)}F(\lambda)\,\dd E_A(\lambda)$。取 $F(\lambda)=\ee^{-|t|\sqrt\lambda}/(2\sqrt\lambda)$，便得到
\[
 G(t)=\frac12A^{-1/2}\ee^{-|t|A^{1/2}}
 =\int_{(0,\infty)}\frac{\ee^{-|t|\sqrt\lambda}}{2\sqrt\lambda}\,\dd E_A(\lambda).
\]
当 $A\ge\mu^2I>0$ 时，$A^{-1/2}$ 有界。对 $t\ne0$，函数演算允许逐谱求导：
\[
 (\partial_t^2-A)G(t)=0.
\]
而一侧导数为 $G'(0^+)=-\frac12I$、$G'(0^-)={+}\frac12I$，故分布二阶导数含 $-\delta(t)I$，从而 $(-\partial_t^2+A)G=\delta(t)I$。


这一计算展示了谱理论在场论中的真正用途：当不能求出所有本征函数时，依然可以用函数演算定义传播子；当可以显式求出谱时，又能够立即得到可计算的级数或积分。谱定理所涉及的定义域、自伴性和连续谱可参照《数学技巧》中自伴算子与谱测度的完整论证。


